\begin{solapartat}
\punts{0,65} $I = \dfrac{Q}{t} \Longrightarrow Q = I\cdot t = 4\times 10^{5}\cdot 5\times 10^{-5} = 20\ \mathrm{C}$

\punts{0,6} $B = \dfrac{\mu_0\cdot I}{2\pi\cdot r} = \dfrac{4\pi\times 10^{-7}\cdot 4\times 10^{5}}{2\pi\cdot 100} = 8\times 10^{-4}\ \mathrm{T}$
\end{solapartat}

\begin{solapartat}
\punts{0,4} $\vec{F} = q\left(\vec{v}\times\vec{B}\right)$

\punts{0,85} Si la velocitat de la partícula és zero, llavors també ho serà la força magnètica; el camp magnètic només actua sobre càrregues en moviment.
\end{solapartat}
